% Auto-generated complete chapter from 20_portfolio_copy_and_site_structure.md
\chapter{20\_portfolio\_copy\_and\_site\_structure}
\label{complete:root_023}

\begin{sourcemeta}
\textbf{Fuente:} \href{file:///E:/Laboral/20_portfolio_copy_and_site_structure.md}{20\_portfolio\_copy\_and\_site\_structure.md}\\
\textbf{Seccion:} root\\
\textbf{Estado:} canonical\\
\textbf{Rol:} Modulo principal\\
\textbf{Lineas originales:} 1246\\
\textbf{Ultima modificacion:} 2026-06-11 15:15\\
\textbf{Deduplicacion conservadora:} 0 bloques repetidos omitidos; 0 caracteres repetidos no reimpresos.
\end{sourcemeta}

\section*{Resumen de orientacion}
This document provides the copy, structure, page architecture and content rules for the public portfolio.

\section*{Contenido completo depurado}
\addcontentsline{toc}{section}{Contenido completo depurado}




\subsection{Purpose}




This document provides the copy, structure, page architecture and content rules for the public portfolio.




It converts the strategy, CV, LinkedIn, GitHub and TwinSight materials into a web-ready structure.




The portfolio should communicate one clear professional identity:




\begin{mdcode}
Bloque de codigo: text
Real-Time 3D Developer / Unity Technical Artist focused on interactive technical visualization, CAD-to-realtime optimization, and Unity WebGL deployment.
\end{mdcode}




The portfolio is not a general personal website. It is a job-search conversion tool.




\medskip\hrule\medskip




\section{1. Portfolio objective}




\subsection{Primary objective}




The portfolio must make a recruiter or technical reviewer understand in under 60 seconds:




\begin{enumerate}
\item what you do;
\item what your strongest project is;
\item why TwinSight is relevant to real jobs;
\item what technologies you can defend;
\item how to contact you;
\item whether you are available for remote work.
\end{enumerate}




\subsection{Secondary objective}




The portfolio must give technical reviewers enough depth to evaluate:




\begin{itemize}
\item Unity/C\# implementation;
\item WebGL deployment;
\item CAD-to-realtime workflow;
\item Blender optimization;
\item runtime interaction systems;
\item technical UI;
\item evaluation method;
\item project limitations;
\item AI-assisted development disclosure.
\end{itemize}




\subsection{What the portfolio should not do}




It should not:




\begin{itemize}
\item look like a generic 3D artist portfolio;
\item lead with FitApp, AI News or old repos;
\item overemphasize full-stack;
\item present ARA as the main project;
\item hide TwinSight behind too many pages;
\item use vague “creative technologist” language without technical proof;
\item rely on long paragraphs;
\item require the viewer to download files;
\item make the user guess what role you want.
\end{itemize}




\medskip\hrule\medskip




\section{2. Recommended portfolio structure}




\subsection{MVP version}




Use this structure first:




\begin{mdcode}
Bloque de codigo: text
Home
Work
  TwinSight X500
  ARA Framework
  Blender Portrait Technical Study
About
Contact
\end{mdcode}




\subsection{Full version}




Use this later:




\begin{mdcode}
Bloque de codigo: text
Home
Work
  TwinSight X500
  ARA Framework
  Blender Portrait Technical Study
  AI News Aggregator [optional]
Technical Notes
About
Contact
\end{mdcode}




\subsection{Do not include initially}




\begin{center}
\scriptsize
\begin{longtable}{>{\raggedright\arraybackslash}p{0.455\linewidth} >{\raggedright\arraybackslash}p{0.455\linewidth}}
\toprule
Item & Reason \\
\midrule
\endfirsthead
\toprule
Item & Reason \\
\midrule
\endhead
FitApp-Free & Low relevance \\
Old web freelance work & Dilutes positioning \\
weak full-stack projects & Lowers technical signal \\
unfinished repos & Credibility risk \\
personal blog & Only if technically relevant \\
unrelated graphic design & Different career track \\
\bottomrule
\end{longtable}
\end{center}




\medskip\hrule\medskip




\section{3. Site navigation}




\subsection{Recommended top navigation}




\begin{mdcode}
Bloque de codigo: text
Work
TwinSight
About
Contact
GitHub
LinkedIn
\end{mdcode}




\subsection{Alternative compact navigation}




\begin{mdcode}
Bloque de codigo: text
Projects
About
Contact
\end{mdcode}




\subsection{Navigation rule}




TwinSight should be accessible in one click from the homepage.




\medskip\hrule\medskip




\section{4. Home page structure}




\subsection{Section 1 — Hero}




\subsubsection{Hero headline}




Recommended:




\begin{mdcode}
Bloque de codigo: text
Real-Time 3D Developer / Unity Technical Artist
\end{mdcode}




\subsubsection{Hero subheadline}




\begin{mdcode}
Bloque de codigo: text
I build interactive technical visualization systems with Unity, WebGL, C#, Blender and CAD-to-realtime optimization.
\end{mdcode}




\subsubsection{Supporting line}




\begin{mdcode}
Bloque de codigo: text
Based in Colombia. Available for remote contractor/B2B roles.
\end{mdcode}




\subsubsection{Primary CTA}




\begin{mdcode}
Bloque de codigo: text
View TwinSight X500
\end{mdcode}




\subsubsection{Secondary CTA}




\begin{mdcode}
Bloque de codigo: text
View GitHub
\end{mdcode}




\subsubsection{Optional third CTA}




\begin{mdcode}
Bloque de codigo: text
Download CV
\end{mdcode}




Only include CV download if the CV is final.




\subsection{Hero layout}




Recommended visual:




\begin{itemize}
\item TwinSight screenshot;
\item drone overview;
\item subtle wireframe overlay;
\item dark/neutral background;
\item no clutter;
\item no generic AI art.
\end{itemize}




Avoid:




\begin{itemize}
\item fantasy/game visuals unless applying to game art;
\item unrelated portrait render as hero;
\item animated background that hurts readability;
\item too much text.
\end{itemize}




\medskip\hrule\medskip




\section{5. Home page copy}




\subsection{Full home page draft}




\begin{mdcode}
Bloque de codigo: text
Real-Time 3D Developer / Unity Technical Artist

I build interactive technical visualization systems with Unity, WebGL, C#, Blender and CAD-to-realtime optimization.

My flagship project, TwinSight X500, is a Unity WebGL prototype for inspecting the assembly of a Holybro X500 V2 drone through component selection, exploded views, cross-section tools, visual modes and technical information panels.

I focus on real-time 3D systems that make complex technical information easier to inspect, understand and communicate.

Based in Colombia. Available for remote contractor/B2B roles.
\end{mdcode}




\subsection{Short version}




\begin{mdcode}
Bloque de codigo: text
I build real-time 3D systems for technical visualization using Unity, C#, WebGL and Blender.

Flagship project: TwinSight X500 — a Unity WebGL drone assembly inspection prototype.
\end{mdcode}




\subsection{More recruiter-friendly version}




\begin{mdcode}
Bloque de codigo: text
Unity / Real-Time 3D developer focused on technical visualization, WebGL deployment and optimized 3D workflows.

Available for remote contractor roles involving Unity, interactive 3D, technical visualization, CAD-to-realtime pipelines, simulation interfaces and WebGL experiences.
\end{mdcode}




\medskip\hrule\medskip




\section{6. Featured projects section}




\subsection{Section title}




\begin{mdcode}
Bloque de codigo: text
Selected Work
\end{mdcode}




\subsection{Intro line}




\begin{mdcode}
Bloque de codigo: text
Projects focused on real-time 3D, technical visualization, Unity WebGL, automation and technical art workflows.
\end{mdcode}




\subsection{Project card order}




\begin{enumerate}
\item TwinSight X500.
\item ARA Framework.
\item Blender Portrait Technical Study.
\item AI News Aggregator only if functional.
\end{enumerate}




\subsection{Project card template}




\begin{mdcode}
Bloque de codigo: text
Project name
One-sentence description
Tags
Primary link
Secondary link
\end{mdcode}




\medskip\hrule\medskip




\section{7. Project card copy}




\subsection{TwinSight X500 card}




\subsubsection{Title}




\begin{mdcode}
Bloque de codigo: text
TwinSight X500
\end{mdcode}




\subsubsection{Subtitle}




\begin{mdcode}
Bloque de codigo: text
Unity WebGL technical visualization prototype for drone assembly inspection.
\end{mdcode}




\subsubsection{Description}




\begin{mdcode}
Bloque de codigo: text
Interactive 3D system for inspecting a Holybro X500 V2 drone assembly using component selection, exploded view, cross-section tools, visual modes, technical UI and CAD-to-realtime optimization.
\end{mdcode}




\subsubsection{Tags}




\begin{mdcode}
Bloque de codigo: text
Unity · C# · WebGL · Blender · CAD-to-Realtime · Technical Visualization
\end{mdcode}




\subsubsection{CTAs}




\begin{mdcode}
Bloque de codigo: text
View Case Study
View GitHub
Watch Demo
\end{mdcode}




\subsection{ARA Framework card}




\subsubsection{Title}




\begin{mdcode}
Bloque de codigo: text
ARA Framework
\end{mdcode}




\subsubsection{Subtitle}




\begin{mdcode}
Bloque de codigo: text
Python research automation prototype.
\end{mdcode}




\subsubsection{Description}




\begin{mdcode}
Bloque de codigo: text
Experimental multi-agent research automation system using Python and LangGraph/LangChain-style workflow orchestration for literature analysis, technical planning and structured report generation.
\end{mdcode}




\subsubsection{Tags}




\begin{mdcode}
Bloque de codigo: text
Python · LangGraph · LangChain · LLM Workflows · Automation · Research Tooling
\end{mdcode}




\subsubsection{CTAs}




\begin{mdcode}
Bloque de codigo: text
View Project
View GitHub
\end{mdcode}




\subsection{Blender portrait card}




\subsubsection{Title}




\begin{mdcode}
Bloque de codigo: text
Hyperrealistic Blender Portrait
\end{mdcode}




\subsubsection{Subtitle}




\begin{mdcode}
Bloque de codigo: text
Technical art study in topology, materials, grooming and lighting.
\end{mdcode}




\subsubsection{Description}




\begin{mdcode}
Bloque de codigo: text
High-fidelity Blender portrait study used as supporting evidence of 3D pipeline literacy, material work, grooming, lighting and rendering fundamentals.
\end{mdcode}




\subsubsection{Tags}




\begin{mdcode}
Bloque de codigo: text
Blender · Topology · Materials · Grooming · Lighting · Rendering
\end{mdcode}




\subsubsection{CTAs}




\begin{mdcode}
Bloque de codigo: text
View Breakdown
ArtStation
\end{mdcode}




\subsection{AI News Aggregator card}




Include only if functional.




\subsubsection{Title}




\begin{mdcode}
Bloque de codigo: text
AI News Aggregator
\end{mdcode}




\subsubsection{Subtitle}




\begin{mdcode}
Bloque de codigo: text
Full-stack AI news aggregation prototype.
\end{mdcode}




\subsubsection{Description}




\begin{mdcode}
Bloque de codigo: text
Experimental FastAPI/React prototype for collecting, organizing and analyzing AI news with a Python backend, database-backed storage and web interface.
\end{mdcode}




\subsubsection{Tags}




\begin{mdcode}
Bloque de codigo: text
Python · FastAPI · React · PostgreSQL · Redis · AI Tooling
\end{mdcode}




\subsubsection{CTAs}




\begin{mdcode}
Bloque de codigo: text
View Project
View GitHub
\end{mdcode}




\medskip\hrule\medskip




\section{8. TwinSight case study page}




\subsection{Page title}




\begin{mdcode}
Bloque de codigo: text
TwinSight X500
\end{mdcode}




\subsection{Subtitle}




\begin{mdcode}
Bloque de codigo: text
Unity WebGL technical visualization prototype for drone assembly inspection.
\end{mdcode}




\subsection{Opening paragraph}




\begin{mdcode}
Bloque de codigo: text
TwinSight X500 is a Unity WebGL interactive technical visualization prototype for inspecting and understanding the assembly of a Holybro X500 V2 drone. The project transforms complex technical documentation and CAD-derived assets into an optimized real-time 3D inspection experience with component selection, exploded views, cross-section tools, visual modes and technical information panels.
\end{mdcode}




\subsection{Core idea block}




\begin{mdcode}
Bloque de codigo: text
Seeing parts is not enough. Users need to understand spatial relationships.
\end{mdcode}




\subsection{Metadata block}




\begin{center}
\scriptsize
\begin{longtable}{>{\raggedright\arraybackslash}p{0.455\linewidth} >{\raggedright\arraybackslash}p{0.455\linewidth}}
\toprule
Field & Value \\
\midrule
\endfirsthead
\toprule
Field & Value \\
\midrule
\endhead
Role & Unity/C\# developer, technical artist, 3D optimization, documentation, evaluation \\
Type & Thesis / portfolio prototype \\
Stack & Unity, C\#, WebGL, URP, UI Toolkit, Shader Graph, Blender \\
Focus & Technical visualization, CAD-to-realtime, interaction systems \\
Status & Thesis complete, pending defense \\
Links & GitHub, demo, video \\
\bottomrule
\end{longtable}
\end{center}




\medskip\hrule\medskip




\section{9. TwinSight case study structure}




\subsection{Recommended sections}




\begin{mdcode}
Bloque de codigo: text
1. Problem
2. Context
3. My role
4. Constraints
5. Solution overview
6. CAD-to-realtime pipeline
7. Interaction systems
8. Visual modes and technical UI
9. WebGL and optimization
10. Evaluation
11. Results
12. Limitations
13. What I would improve
14. AI-assisted development note
15. Asset/manufacturer disclaimer
\end{mdcode}




\medskip\hrule\medskip




\section{10. TwinSight case study copy}




\subsection{Problem}




\begin{mdcode}
Bloque de codigo: text
Drone assembly documentation often relies on 2D diagrams, manuals and static references. These materials can show individual parts, but they require users to mentally reconstruct spatial relationships between components, fasteners, subassemblies and assembly steps.

TwinSight X500 explores whether an interactive real-time 3D viewer can make those spatial relationships easier to inspect and understand.
\end{mdcode}




\subsection{Context}




\begin{mdcode}
Bloque de codigo: text
The project was developed as a Multimedia Engineering thesis. It uses the Holybro X500 V2 drone as a technical assembly case and focuses on converting complex technical documentation and CAD-derived assets into a browser-accessible Unity WebGL experience.
\end{mdcode}




\subsection{My role}




\begin{mdcode}
Bloque de codigo: text
I handled Unity/C# implementation, runtime interaction systems, technical UI, CAD-to-realtime asset optimization, Blender-based asset preparation, documentation, evaluation design and final integration.

AI-assisted tools were used as support for implementation, debugging, research and documentation. Architecture decisions, integration, testing, validation and final technical ownership remained my responsibility.
\end{mdcode}




\subsection{Constraints}




\begin{mdcode}
Bloque de codigo: text
The project had to balance technical clarity, real-time interaction, WebGL deployment constraints, academic scope, asset optimization, usability evaluation and limited production time.

The goal was not to build a full engineering simulation. The goal was to create an interactive technical visualization prototype for assembly inspection and spatial understanding.
\end{mdcode}




\subsection{Solution overview}




\begin{mdcode}
Bloque de codigo: text
TwinSight X500 provides an interactive 3D inspection environment where users can select components, view technical information, explore exploded views, use cross-section tools and switch between visual modes such as realistic, ghosted, blueprint, wireframe and qualitative thermal-style visualization.
\end{mdcode}




\subsection{CAD-to-realtime pipeline}




\begin{mdcode}
Bloque de codigo: text
The source geometry came from CAD/manufacturer-derived assets. The pipeline involved conversion, inspection, cleanup, optimization, Blender-based preparation, UV work, material planning, Unity integration and WebGL-conscious asset handling.

The goal was to reduce heavy CAD geometry while preserving component identity and assembly relationships.
\end{mdcode}




\subsection{Optimization copy}




\begin{mdcode}
Bloque de codigo: text
The project reduced dense CAD-derived geometry from over 6.5 million triangles to approximately 95,617 optimized triangles. [verify final metric]

This reduction was necessary to make the model suitable for real-time browser-based visualization while preserving the visual information needed for technical inspection.
\end{mdcode}




\subsection{Interaction systems}




\begin{mdcode}
Bloque de codigo: text
The runtime systems were designed around inspection rather than passive viewing.

Core interactions include:

- component selection;
- highlighting;
- technical information panels;
- exploded view;
- cross-section / clipping;
- visual mode switching;
- mobile/web-oriented UI behavior.
\end{mdcode}




\subsection{Visual modes}




\begin{mdcode}
Bloque de codigo: text
The visual modes allow users to inspect the drone from different technical perspectives. Realistic mode supports general recognition, wireframe and blueprint modes emphasize structure, ghosted/X-ray modes help reveal hidden relationships, and qualitative thermal-style visualization explores an alternative analytical layer.
\end{mdcode}




\subsection{Evaluation}




\begin{mdcode}
Bloque de codigo: text
The project was evaluated using usability and workload methods including SUS, NASA-TLX Raw and Think-Aloud methodology. [verify exact final metrics]

The evaluation was designed to compare interactive 3D support with conventional 2D reference material in assembly-inspection tasks.
\end{mdcode}




\subsection{Results}




Use only after verifying final numbers:




\begin{mdcode}
Bloque de codigo: text
Reported evaluation data included:

- 12 validation participants;
- SUS average: 91.88;
- NASA-TLX Raw: 8.69 for the 3D viewer;
- NASA-TLX Raw: 19.89 for the 2D support condition;
- 96 task-condition records.

These results should be presented as academic prototype findings, not as commercial product validation.
\end{mdcode}




\subsection{Limitations}




\begin{mdcode}
Bloque de codigo: text
TwinSight X500 is an academic prototype, not a commercial production product. It does not replace engineering CAD software, formal simulation tools, FEA analysis or manufacturer documentation.

Some visual modes are qualitative. Some asset usage may require public-display limitations. The public portfolio version may not include every model file from the original thesis project.
\end{mdcode}




\subsection{What I would improve}




\begin{mdcode}
Bloque de codigo: text
Given more time, I would improve build profiling, formalize LOD handling, refine mobile interaction details, add more structured user analytics, improve the public WebGL deployment pipeline, expand technical documentation and separate the prototype into clearer runtime modules.
\end{mdcode}




\subsection{AI-assisted development note}




\begin{mdcode}
Bloque de codigo: text
AI-assisted workflows were used during parts of development for implementation support, debugging, research and documentation. AI was not used as a replacement for technical ownership.

I remained responsible for architecture decisions, integration, validation, testing, tradeoff analysis and final project delivery.
\end{mdcode}




\subsection{Asset/manufacturer disclaimer}




\begin{mdcode}
Bloque de codigo: text
This project is an academic/portfolio prototype created for technical visualization and assembly-understanding purposes.

References to the Holybro X500 V2 platform and any CAD/manufacturer-derived assets are used for educational and technical visualization purposes only.

This project is not affiliated with, endorsed by, sponsored by, or officially supported by Holybro.
\end{mdcode}




\medskip\hrule\medskip




\section{11. ARA Framework project page}




\subsection{Page title}




\begin{mdcode}
Bloque de codigo: text
ARA Framework
\end{mdcode}




\subsection{Subtitle}




\begin{mdcode}
Bloque de codigo: text
AI-assisted research automation prototype.
\end{mdcode}




\subsection{Opening paragraph}




\begin{mdcode}
Bloque de codigo: text
ARA Framework is a Python-based research automation prototype exploring multi-agent workflows for niche analysis, literature research, technical planning, implementation guidance and structured report generation.
\end{mdcode}




\subsection{Positioning note}




\begin{mdcode}
Bloque de codigo: text
ARA is a secondary project. It supports the profile by showing automation, tooling, workflow architecture and technical documentation. It is not positioned as a production AI platform or as the main career direction.
\end{mdcode}




\subsection{Metadata}




\begin{center}
\scriptsize
\begin{longtable}{>{\raggedright\arraybackslash}p{0.455\linewidth} >{\raggedright\arraybackslash}p{0.455\linewidth}}
\toprule
Field & Value \\
\midrule
\endfirsthead
\toprule
Field & Value \\
\midrule
\endhead
Type & Prototype \\
Stack & Python, LangGraph, LangChain, Redis, APIs \\
Focus & Research automation, LLM workflows, technical documentation \\
Status & Experimental / prototype \\
Role & Architecture, implementation, documentation \\
\bottomrule
\end{longtable}
\end{center}




\subsection{Case study sections}




\begin{mdcode}
Bloque de codigo: text
1. Problem
2. Workflow concept
3. Architecture
4. Tools
5. What it demonstrates
6. Limitations
7. Next steps
\end{mdcode}




\subsection{Copy}




\begin{mdcode}
Bloque de codigo: text
The goal of ARA Framework is to test how LLM-assisted workflows can structure parts of academic and technical research. The system explores modular agent-style stages such as niche analysis, literature research, technical architecture planning, implementation guidance and report synthesis.

The project demonstrates Python automation, workflow architecture, structured output design and technical documentation. It does not claim to replace expert review or operate as a production research system.
\end{mdcode}




\medskip\hrule\medskip




\section{12. Blender portrait project page}




\subsection{Page title}




\begin{mdcode}
Bloque de codigo: text
Hyperrealistic Blender Portrait
\end{mdcode}




\subsection{Subtitle}




\begin{mdcode}
Bloque de codigo: text
Technical art study in topology, materials, grooming and lighting.
\end{mdcode}




\subsection{Opening paragraph}




\begin{mdcode}
Bloque de codigo: text
This project is a high-fidelity Blender portrait study focused on realistic character pipeline fundamentals, including topology, material definition, grooming, lighting and rendering.
\end{mdcode}




\subsection{Positioning note}




\begin{mdcode}
Bloque de codigo: text
This project supports the technical art side of the portfolio. It is not the main career-positioning project and should not displace TwinSight X500.
\end{mdcode}




\subsection{Sections}




\begin{mdcode}
Bloque de codigo: text
1. Goal
2. Process
3. Topology
4. Materials
5. Grooming
6. Lighting
7. Final render
8. Lessons learned
\end{mdcode}




\subsection{Copy}




\begin{mdcode}
Bloque de codigo: text
The project was made as a structured technical art study following non-certified training material. It demonstrates 3D pipeline literacy, attention to surface detail, topology awareness, material work and rendering fundamentals.

Its role in the portfolio is to support the technical artist positioning, especially when applying to Unity Technical Artist or real-time 3D roles that value asset-pipeline awareness.
\end{mdcode}




\medskip\hrule\medskip




\section{13. About page}




\subsection{About headline}




\begin{mdcode}
Bloque de codigo: text
I build interactive real-time 3D systems for technical visualization.
\end{mdcode}




\subsection{About body}




\begin{mdcode}
Bloque de codigo: text
I am a Colombia-based Real-Time 3D Developer / Unity Technical Artist focused on Unity WebGL, C#, CAD-to-realtime optimization and interactive technical visualization.

My strongest work is TwinSight X500, a Unity WebGL technical visualization prototype for inspecting drone assembly. The project combines real-time interaction systems, optimized CAD-derived assets, technical UI, visual modes and usability/workload evaluation.

My background combines Multimedia Engineering, advanced Electronic Engineering coursework, computer graphics, Blender-based asset workflows, Python-assisted tooling and technical documentation.

I am interested in roles involving:

- Unity development;
- real-time 3D;
- Unity WebGL;
- technical visualization;
- digital twin visualization;
- simulation interfaces;
- XR-adjacent development;
- technical art workflows;
- Python automation for production or research tooling.

I am based in Colombia and available for remote contractor/B2B work.
\end{mdcode}




\subsection{Short About version}




\begin{mdcode}
Bloque de codigo: text
Colombia-based Real-Time 3D Developer / Unity Technical Artist focused on Unity WebGL, C#, CAD-to-realtime optimization and interactive technical visualization.

Main project: TwinSight X500, a Unity WebGL drone assembly inspection prototype.

Available for remote contractor/B2B roles.
\end{mdcode}




\medskip\hrule\medskip




\section{14. Contact page}




\subsection{Contact headline}




\begin{mdcode}
Bloque de codigo: text
Available for remote Unity, real-time 3D and technical visualization work.
\end{mdcode}




\subsection{Contact copy}




\begin{mdcode}
Bloque de codigo: text
I am currently based in Colombia and available for remote contractor/B2B roles involving Unity, real-time 3D, Unity WebGL, technical visualization, CAD-to-realtime workflows, interactive 3D systems and Python-assisted tooling.

For relevant opportunities, contact me through email or LinkedIn.
\end{mdcode}




\subsection{Contact links}




\begin{mdcode}
Bloque de codigo: text
Email: [email]
LinkedIn: [LinkedIn URL]
GitHub: [GitHub URL]
ArtStation: [ArtStation URL]
CV: [CV URL]
\end{mdcode}




\subsection{Optional work authorization note}




\begin{mdcode}
Bloque de codigo: text
Current availability: remote contractor/B2B from Colombia.
Medium-term EU mobility may become relevant through expected Portuguese citizenship/passport, but current applications should be treated as Colombia-based remote availability.
\end{mdcode}




Do not over-detail legal pathways on the contact page unless necessary.




\medskip\hrule\medskip




\section{15. Footer copy}




\begin{mdcode}
Bloque de codigo: text
© Alexander Woodcock Salomón
Real-Time 3D · Unity WebGL · Technical Visualization · CAD-to-Realtime
\end{mdcode}




Alternative:




\begin{mdcode}
Bloque de codigo: text
Built for technical visualization, real-time 3D and Unity WebGL work.
\end{mdcode}




\medskip\hrule\medskip




\section{16. Visual direction}




\subsection{Recommended style}




\begin{center}
\scriptsize
\begin{longtable}{>{\raggedright\arraybackslash}p{0.455\linewidth} >{\raggedright\arraybackslash}p{0.455\linewidth}}
\toprule
Element & Direction \\
\midrule
\endfirsthead
\toprule
Element & Direction \\
\midrule
\endhead
Theme & Dark or neutral \\
Visual language & Technical, clean, precise \\
Main visual & TwinSight screenshot \\
Accent & Wireframe / CAD lines \\
Typography & Clean sans-serif \\
Animation & Minimal \\
Layout & Fast scanning \\
Tone & Technical, direct \\
\bottomrule
\end{longtable}
\end{center}




\subsection{Avoid}




\begin{itemize}
\item heavy decorative animations;
\item visual noise;
\item huge paragraphs;
\item generic AI-generated backgrounds;
\item too much gaming/fantasy art;
\item unrelated profile photo as hero;
\item over-stylized “creative portfolio” look.
\end{itemize}




\subsection{Portfolio should feel like}




\begin{mdcode}
Bloque de codigo: text
technical product case study + real-time 3D portfolio
\end{mdcode}




not:




\begin{mdcode}
Bloque de codigo: text
graphic design portfolio
\end{mdcode}




\medskip\hrule\medskip




\section{17. Case study media checklist}




\subsection{TwinSight media}




Required:




\begin{itemize}
\item hero screenshot;
\item 60–120 second demo video;
\item component selection screenshot;
\item exploded view screenshot;
\item cross-section screenshot;
\item visual modes screenshot;
\item CAD-to-realtime before/after image;
\item optimization table;
\item UI screenshot;
\item evaluation summary image/table.
\end{itemize}




Optional:




\begin{itemize}
\item short GIF loop;
\item architecture diagram;
\item pipeline diagram;
\item mobile layout screenshot;
\item code snippet screenshot.
\end{itemize}




\subsection{ARA media}




Required if published:




\begin{itemize}
\item architecture diagram;
\item workflow screenshot;
\item sample output screenshot;
\item README screenshot or terminal screenshot.
\end{itemize}




Optional:




\begin{itemize}
\item short CLI demo;
\item report generation preview.
\end{itemize}




\subsection{Blender media}




Required:




\begin{itemize}
\item final render;
\item topology screenshot;
\item material/lighting breakdown;
\item grooming close-up;
\item progression image.
\end{itemize}




\medskip\hrule\medskip




\section{18. Site copy consistency rules}




Use the same phrasing across portfolio, CV, LinkedIn and GitHub.




\subsection{Preferred terms}




Use:




\begin{mdcode}
Bloque de codigo: text
Real-Time 3D
Unity WebGL
Technical Visualization
CAD-to-Realtime
Interactive 3D
Technical UI
Runtime Interaction Systems
Asset Optimization
Blender-based preparation
\end{mdcode}




Avoid overusing:




\begin{mdcode}
Bloque de codigo: text
multimedia
creative
generalist
full-stack
AI expert
game developer
graphic design
\end{mdcode}




\subsection{Thesis wording}




Correct:




\begin{mdcode}
Bloque de codigo: text
Thesis project completed, pending defense.
\end{mdcode}




Avoid:




\begin{mdcode}
Bloque de codigo: text
Currently starting thesis development.
\end{mdcode}




\subsection{AI wording}




Correct:




\begin{mdcode}
Bloque de codigo: text
AI-assisted workflows were used as implementation, debugging, research and documentation support.
\end{mdcode}




Avoid:




\begin{mdcode}
Bloque de codigo: text
AI built this project.
\end{mdcode}




\medskip\hrule\medskip




\section{19. SEO and recruiter keywords}




Use naturally in page headings, project tags and descriptions.




\subsection{Core keywords}




\begin{mdcode}
Bloque de codigo: text
Unity Developer
Unity Technical Artist
Real-Time 3D Developer
Unity WebGL Developer
Technical Visualization Developer
Interactive 3D Developer
CAD-to-Realtime
Digital Twin Visualization
Simulation Visualization
3D Optimization
Blender
C#
WebGL
Technical UI
\end{mdcode}




\subsection{Secondary keywords}




\begin{mdcode}
Bloque de codigo: text
Python Automation
LLM Workflows
Research Automation
Shader Graph
URP
UI Toolkit
XR
AR/VR
Product Visualization
Assembly Visualization
\end{mdcode}




\subsection{Do not keyword stuff}




The text must remain readable.




\medskip\hrule\medskip




\section{20. Portfolio MVP build options}




\subsection{Acceptable MVP platforms}




\begin{center}
\scriptsize
\begin{longtable}{>{\raggedright\arraybackslash}p{0.455\linewidth} >{\raggedright\arraybackslash}p{0.455\linewidth}}
\toprule
Platform & Use case \\
\midrule
\endfirsthead
\toprule
Platform & Use case \\
\midrule
\endhead
GitHub Pages & Developer-facing, free \\
Framer & Fast visual site \\
Webflow & Polished no-code site \\
Notion / Super & Very fast MVP \\
Custom HTML/CSS & Full control \\
WordPress & Only if already fast to deploy \\
\bottomrule
\end{longtable}
\end{center}




\subsection{Recommendation}




Use the fastest option that allows:




\begin{itemize}
\item case study page;
\item project cards;
\item screenshots;
\item video embeds;
\item external links;
\item readable design;
\item easy updates.
\end{itemize}




Do not spend weeks building the portfolio tech stack.




The portfolio content is more important than the framework.




\medskip\hrule\medskip




\section{21. Portfolio MVP acceptance criteria}




The MVP is ready to use when:




\begin{itemize}
\item homepage clearly states role;
\item TwinSight appears first;
\item TwinSight case study is readable;
\item at least 6 TwinSight screenshots exist;
\item GitHub and LinkedIn links work;
\item contact email works;
\item no false claims;
\item no outdated thesis status;
\item mobile layout is acceptable;
\item page loads reasonably fast;
\item CV link works if included.
\end{itemize}




\subsection{Not required for MVP}




\begin{itemize}
\item perfect animations;
\item full blog;
\item all projects complete;
\item custom CMS;
\item advanced scroll effects;
\item perfect bilingual version;
\item 3D hero embed;
\item multiple case studies.
\end{itemize}




\medskip\hrule\medskip




\section{22. Portfolio full version backlog}




Add later:




\begin{itemize}
\item bilingual English/Spanish toggle;
\item technical notes section;
\item WebGL embedded demo;
\item more detailed pipeline diagrams;
\item downloadable case study PDF;
\item project filtering;
\item ARA demo;
\item updated Blender breakdown;
\item blog posts on Unity WebGL optimization.
\end{itemize}




\medskip\hrule\medskip




\section{23. Suggested portfolio file structure}




If building custom/static:




\begin{mdcode}
Bloque de codigo: text
portfolio/
  index.html
  about.html
  contact.html
  work/
    twinsight-x500.html
    ara-framework.html
    blender-portrait.html
  assets/
    images/
      twinsight/
      ara/
      blender/
    video/
    css/
    js/
  docs/
    cv.pdf
\end{mdcode}




If using React:




\begin{mdcode}
Bloque de codigo: text
src/
  pages/
    Home.jsx
    About.jsx
    Contact.jsx
    projects/
      TwinSight.jsx
      ARA.jsx
      BlenderPortrait.jsx
  components/
    ProjectCard.jsx
    Hero.jsx
    Section.jsx
  assets/
\end{mdcode}




Do not overbuild unless there is time.




\medskip\hrule\medskip




\section{24. Portfolio review checklist}




Before sharing publicly:




\begin{itemize}
\item [ ] No typos.
\item [ ] English is clean.
\item [ ] TwinSight metrics verified.
\item [ ] Links work.
\item [ ] Images load.
\item [ ] Mobile view works.
\item [ ] Contact works.
\item [ ] CV link works.
\item [ ] GitHub is not embarrassing.
\item [ ] No old thesis status.
\item [ ] ARA not overclaimed.
\item [ ] AI disclosure present.
\item [ ] Asset disclaimer present.
\item [ ] Site communicates target role in 5 seconds.
\end{itemize}




\medskip\hrule\medskip




\section{25. Final homepage draft}




Use this as the first implementation copy.




\begin{mdcode}
Bloque de codigo: text
Real-Time 3D Developer / Unity Technical Artist

I build interactive technical visualization systems with Unity, WebGL, C#, Blender and CAD-to-realtime optimization.

My flagship project, TwinSight X500, is a Unity WebGL prototype for inspecting the assembly of a Holybro X500 V2 drone through component selection, exploded views, cross-section tools, visual modes and technical information panels.

I focus on real-time 3D systems that make complex technical information easier to inspect, understand and communicate.

Based in Colombia. Available for remote contractor/B2B roles.

[View TwinSight X500] [View GitHub] [Contact]
\end{mdcode}




\medskip\hrule\medskip




\section{26. Final Work page draft}




\begin{mdcode}
Bloque de codigo: text
Selected Work

Projects focused on real-time 3D, technical visualization, Unity WebGL, automation and technical art workflows.

TwinSight X500
Unity WebGL technical visualization prototype for drone assembly inspection.
Unity · C# · WebGL · Blender · CAD-to-Realtime · Technical Visualization
[View Case Study]

ARA Framework
Python research automation prototype using multi-agent workflow orchestration for literature analysis, technical planning and structured report generation.
Python · LangGraph · LangChain · LLM Workflows · Automation
[View Project]

Hyperrealistic Blender Portrait
Technical art study in topology, materials, grooming, lighting and rendering.
Blender · Topology · Materials · Grooming · Lighting
[View Breakdown]
\end{mdcode}




\medskip\hrule\medskip




\section{27. Final About page draft}




\begin{mdcode}
Bloque de codigo: text
I build interactive real-time 3D systems for technical visualization.

I am a Colombia-based Real-Time 3D Developer / Unity Technical Artist focused on Unity WebGL, C#, CAD-to-realtime optimization and interactive technical visualization.

My strongest work is TwinSight X500, a Unity WebGL technical visualization prototype for inspecting drone assembly. The project combines real-time interaction systems, optimized CAD-derived assets, technical UI, visual modes and usability/workload evaluation.

My background combines Multimedia Engineering, advanced Electronic Engineering coursework, computer graphics, Blender-based asset workflows, Python-assisted tooling and technical documentation.

I am interested in roles involving Unity development, real-time 3D, Unity WebGL, technical visualization, digital twin visualization, simulation interfaces, XR-adjacent development, technical art workflows and Python automation for production or research tooling.

I am based in Colombia and available for remote contractor/B2B work.
\end{mdcode}




\medskip\hrule\medskip




\section{28. Final Contact page draft}




\begin{mdcode}
Bloque de codigo: text
Available for remote Unity, real-time 3D and technical visualization work.

I am currently based in Colombia and available for remote contractor/B2B roles involving Unity, real-time 3D, Unity WebGL, technical visualization, CAD-to-realtime workflows, interactive 3D systems and Python-assisted tooling.

Contact:
Email: [email]
LinkedIn: [LinkedIn URL]
GitHub: [GitHub URL]
ArtStation: [ArtStation URL]
CV: [CV URL]
\end{mdcode}




\medskip\hrule\medskip




\section{29. Next file dependency}




The next file should be:




\begin{mdcode}
Bloque de codigo: text
21_twinsight_demo_video_script_and_shotlist.md
\end{mdcode}




It should define:




\begin{itemize}
\item 60-second demo script;
\item 90-second demo script;
\item 120-second demo script;
\item shot list;
\item capture checklist;
\item captions;
\item voiceover;
\item LinkedIn cutdown;
\item silent version;
\item technical breakdown version.
\end{itemize}




No Deep Research or agent required.



